% authority: science_draft
% claim_status: draft/control
\documentclass[11pt]{article}
\usepackage[margin=1in]{geometry}
\usepackage{amsmath,amssymb,amsthm,booktabs,array,longtable,enumitem,hyperref}
\hypersetup{hidelinks}
\setlist{nosep}

\title{V22 P4-T02 B2 Descriptor Formalization v1\\
\large Typed source-extension law and two conditional results}
\author{Ontology Formalizer\\\small Internal parent-child synthesis}
\date{2026-08-09}

\theoremstyle{definition}
\newtheorem{definition}{Definition}
\theoremstyle{plain}
\newtheorem{proposition}{Proposition}
\newtheorem{theorem}{Theorem}
\newtheorem{corollary}{Corollary}
\newcommand{\A}{\mathcal{A}}
\newcommand{\Sectors}{\mathcal{S}}
\newcommand{\Desc}{\mathcal{D}_{\mathrm{B2}}}
\newcommand{\Psrc}{\mathcal{P}_{\mathrm{src}}}
\newcommand{\Proj}{\mathbb{P}}
\newcommand{\Char}{\mathcal{C}}
\newcommand{\J}{\mathrm{J}}
\setlength{\emergencystretch}{3em}

\begin{document}
\raggedbottom
\maketitle

\begin{center}
\fbox{\parbox{0.92\textwidth}{
\textbf{Status and authority.} This is a task-local \emph{draft/control},
\emph{proposal-only} item of \emph{source-extension data}. It formalizes a
descriptor schema and proves conditional mathematical statements about that
schema. It does not instantiate a B2 continuum lift, adopt a source law, modify
canonical ontology, activate B2, satisfy P2-T01, unlock P4-T03, derive a
physical cone or metric, issue a Gate verdict, or authorize promotion or
external action.
}}
\end{center}

\begin{abstract}
The inactive B2 fallback retains the exact finite P7 matter package but lacks
a typed continuum lift and sector-common principal law. We define the complete
ten-component descriptor envelope selected by RT-20260809-020, state seven
fail-closed proof-obligation families, prove a projective common-principal
gluing theorem that assumes no quadratic or Lorentzian target, and prove that
finite sampled traces cannot by themselves determine a continuum principal
class. The specification is complete; no descriptor instance is constructed.
Consequently the B2 activation-readiness vector remains
$(1,1,1,1,0,0)$ and P4-T03 remains locked.
\end{abstract}

\section{Fixed boundary and missing source law}

Let $\mathfrak P_7$ denote the exact protected finite P7 conditional input,
including its declared sector, transition, device, propagation-profile,
coupling, variational, and operational records. Its protected interpretation
is adopted by explicit human postulate in its finite scope; it is not a
continuum matter theory and supplies no source-to-target constructor. We do
not change any member or hash of $\mathfrak P_7$.

The candidate family is
\nolinkurl{FAM-V22-B2-MATTER-PRINCIPAL-POLYNOMIAL}; its exact candidate identity
is \nolinkurl{CAND-V22-B2-P7-COMMON-PRINCIPAL-LIFT-V1}. The current ontology does
not derive a target-atlas-free finite-P7-to-continuum lift, continuum sector
operators, or a source compatibility law selecting one common reduced
principal target. This is the
\nolinkurl{derivation_critical_missing_source_law} trigger. It does not say
that such a law is impossible.

The source arena $\A$ below is the already declared smooth four-dimensional
source arena or a proposal-only open patch thereof. It is not physical
spacetime. Source charts are computational presentations and are not target
atlases.

\section{Typed descriptor law}

\begin{definition}[B2 descriptor envelope]
A B2 continuum-lift descriptor instance over a source domain $U\subset\A$ is
the tuple
\begin{equation}
\Desc=
(\Sectors,\{F_s\}_{s\in\Sectors},L,Q_{\rm src},Q_{\rm out},
\{E_s\}_{s\in\Sectors},C,\Psrc,\Omega,N).
\label{eq:descriptor}
\end{equation}
Its components have the following types.
\begin{enumerate}[label=\textbf{D\arabic*.}]
  \item $\Sectors$ is exactly the finite equipped-sector index set exposed by
  $\mathfrak P_7$; no Standard Model or target matter labels may be added.
  \item Each $F_s\to U$ is a finite-rank real source bundle with a declared
  regularity class $C^{r_s}$, $r_s\ge m_s+1$, and a declared regular locus.
  The lift $L=(L_{s,n})$ is a natural family
  \[
    L_{s,n}:\mathfrak P_{7,s}|_{K_n}\longrightarrow
    \Gamma^{r_s}(U,F_s)
  \]
  over a directed source-refinement system $(K_n,\rho_{n+1,n})$, with
  sampling maps $S_{s,n}$ and an explicit consistency norm or topology.
  \item $Q_{\rm src}$ and $Q_{\rm out}$ declare the source and output
  equivalence groupoids, their regular quotient tangents, singular strata,
  and the descent maps used by the P2-T01 interface. Passive source-coordinate
  change is representation covariance, not automatically a physical quotient.
  \item Each sector equation is a local source operator
  \[
    E_s:\J^{m_s}(F_s)\longrightarrow W_s
  \]
  with constraints, vertical redundancies, equation redundancies, a regular
  background locus, and the reduced linearized symbol
  $\bar\sigma_s(x,k):V^{\rm red}_{s,x}\to W^{\rm red}_{s,x}$ declared.
  \item $C$ is a source-defined compatibility relation on the projective
  reduced characteristic hypersurfaces of the sectors. It may use only the
  objects above and source-operational data in $\mathfrak P_7$; it may not
  select a desired GR cone or target geometry.
  \item $\Psrc$ is a common projective principal class. Locally it is
  represented by a reduced homogeneous polynomial $P_i(x,k)$, defined only
  up to a nowhere-zero source unit. Its zero set must be derived from every
  sector's reduced symbol through $C$. No degree, quadratic factor,
  Lorentzian signature, time orientation, or physical scale is assumed.
  \item $\Omega$ is a source-to-operational typing and no-target-import
  receipt. $N$ is a total procedure that maps a completed descriptor instance
  into exactly one allowed P2-T01 verdict while preserving the rule that
  passage is necessary but not sufficient.
\end{enumerate}
\end{definition}

The tuple in \eqref{eq:descriptor} is a \emph{schema} in this packet. A B2
descriptor instance exists only after every component is supplied and every
blocking obligation below is evidenced on the same declared domain.

\section{Reduced sector symbols and the common target}

At a regular point $x\in U$, let
$p_{s,x}(k)=\det\bar\sigma_s(x,k)$ when the reduced symbol is square. More
generally, the descriptor must name the invariant characteristic ideal used
when a determinant is not lawful. Write
\[
 \Char_{s,x}=Z_{\Proj(T_x^*\A)}(\operatorname{rad}p_{s,x})
\]
for the reduced projective characteristic hypersurface in the determinant
case. It is a source-cotangent object, not yet a physical causal cone.

\begin{definition}[Sector-common projective principal class]
The relation $C$ succeeds on a regular patch precisely when there is a reduced
projective hypersurface $\Char_x\subset\Proj(T_x^*\A)$ such that
\begin{equation}
  \Char_{s,x}=\Char_x\quad\text{for every declared }s\in\Sectors,
  \label{eq:commonzero}
\end{equation}
and the local defining polynomials transform by nowhere-zero source units on
overlaps. Any local generator $P_i$ represents $\Psrc=[P]$. Equality of one
chosen coefficient normalization is not required; equality of the reduced
projective characteristic set and coherent unit class is required.
\end{definition}

This definition permits higher-degree and non-quadratic candidates. It also
fails closed when the declared sectors have inequivalent characteristic sets,
when a common set is chosen only by discarding a sector, or when the only
selector is resemblance to the GR target.

% THM-V22-P4T02-B2-PROJECTIVE-COMMON-PRINCIPAL-GLUING
\begin{theorem}[Projective common-principal gluing]
Let $\{U_i\}$ cover a regular source domain $U$. Suppose each $U_i$ carries a
reduced homogeneous polynomial $P_i$ whose projective zero set equals every
declared sector characteristic set on $U_i$. Suppose further that on every
nonempty overlap
\[
  P_i=g_{ij}P_j,
  \qquad g_{ij}\in C^{r}(U_i\cap U_j)^{\times},
\]
with $g_{ii}=1$, $g_{ij}g_{ji}=1$, and
$g_{ij}g_{jk}g_{ki}=1$ on triple overlaps. Then the local zero sets glue to a
unique global projective characteristic hypersurface $\Char\subset
\Proj(T^*U)$, and the $P_i$ define a section of a source line bundle whose
projective class $\Psrc$ is independent of source-chart and local-frame
choice. Any other family with the same reduced local hypersurfaces differs
locally by nowhere-zero units, subject to the same cocycle law.
\end{theorem}

\begin{proof}
Multiplication by a nowhere-zero unit does not change a polynomial's
projective zero set, so the local sets agree on pairwise overlaps and glue as
a unique subset $\Char$. The functions $g_{ij}$ satisfy the transition laws
for a line bundle; the $P_i$ therefore form a section of that bundle. Changing
a source chart or field/equation frame changes a representative by the
declared transition unit but not its projective class. The final statement is
the same overlap argument for the ratio of two reduced local generators.
No quadratic form, signature, target atlas, or benchmark behavior enters the
construction.
\end{proof}

\paragraph{Exact limitation.}
The theorem is conditional. It proves what coherent local common-principal
data would mean and how it would descend; it does not prove that the P7 package
supplies the $F_s$, $E_s$, $C$, $P_i$, or cocycle.

\section{Finite-trace underdetermination}

The finite P7 input cannot be turned into a continuum operator merely by
renaming finite sectors as fields. The following theorem isolates the missing
law without asserting a global no-go.

% THM-V22-P4T02-B2-FINITE-TRACE-PRINCIPAL-UNDERDETERMINATION
\begin{theorem}[Finite traces do not determine the continuum principal class]
Let $\Xi\subset U$ be a finite source sampling set and let $r\ge0$. Suppose a
proposed lift constrains continuum data only through $r$-jets on $\Xi$. For
any local source operator family $E_s$ of order $m$ and any order-$m$ source
operator family $B_s$, there exists a nonzero smooth source scalar $f$ with
$j^r_{\Xi}f=0$ such that
\[
  \widetilde E_s=E_s+fB_s
\]
has the same admitted finite $r$-jet traces on $\Xi$ but may have a different
principal symbol on an open subset of $U\setminus\Xi$. In particular, choosing
all $E_s=\partial_0$ and
$B_s=c_s\partial_1$ with distinct real $c_s$ gives one completion with the
common representative $k_0$ and another, away from $f=0$, with distinct
sector representatives $k_0+fc_sk_1$. Hence finite sampled traces alone do not
select whether \eqref{eq:commonzero} holds.
\end{theorem}

\begin{proof}
Because $\Xi$ is finite, choose pairwise source-coordinate neighborhoods and
a smooth function that vanishes through order $r$ at every point of $\Xi$ but
is nonzero on some open subset. Leibniz expansion shows that every admitted
trace through order $r$ of $fB_s$ vanishes on $\Xi$, so $E_s$ and
$\widetilde E_s$ agree on the declared finite traces. Their principal symbols
differ by $f\sigma(B_s)$ wherever $f\ne0$. For the explicit first-order
families, the unperturbed characteristic hyperplane is $k_0=0$ for every
sector, whereas the perturbed hyperplanes $k_0+fc_sk_1=0$ are pairwise
distinct when $f\ne0$ and the $c_s$ are distinct. Thus the finite trace cannot
choose the common-principal property.
\end{proof}

\begin{corollary}[Missing selector law]
Any successful B2 instance must add and justify at least one source-side
condition stronger than finite trace agreement: for example jet-density,
operator naturality, a source variational construction, refinement stability,
or an independently motivated compatibility selector. The current ontology
does not derive which condition is correct.
\end{corollary}

This is a scoped underdetermination result for finite sampled interfaces. It
does not exclude nonlocal lifts, dense refinement limits, or future source
extensions, and it does not refute B2 globally.

\section{Seven blocking proof-obligation families}

\begin{longtable}{>{\raggedright\arraybackslash}p{0.09\textwidth}>{\raggedright\arraybackslash}p{0.29\textwidth}>{\raggedright\arraybackslash}p{0.54\textwidth}}
\toprule
ID & Object & Fail-closed obligation\\
\midrule
D1 & fields and regularity & Declare every bundle, rank, source domain, differentiability order, regular stratum, singular stratum, and permitted variation.\\
D2 & finite-to-continuum lift & Prove typed naturality, sampling consistency, refinement or limit control, finite-data recovery, and a non-circular selector stronger than finite trace agreement.\\
D3 & quotients & Declare source and output equivalences, constant-rank loci, invariant observables, and representative-independent descent.\\
D4 & sector equations & Construct $E_s$, constraints, redundancies, reduced perturbation spaces, and covariant principal symbols before taking determinants.\\
D5 & compatibility and $\Psrc$ & Derive $C$, prove sector-common projective equality, unit/cocycle coherence, and finite-variation robustness without imposing a desired cone.\\
D6 & operational and no-target typing & Bind every response to source type, device bridge, scale status, domain, invariant readout, and falsifier; reject target atlas, metric, tetrad, coframe, benchmark fit, or validator authority as premises.\\
D7 & adequacy interface & Populate all ten P2-T01 fields and emit exactly one allowed necessary-condition verdict; never infer sufficiency or Gate B passage.\\
\bottomrule
\end{longtable}

The task-local proof-obligation matrix records 35 atomic checks under these
seven families. Every family is blocking for an instantiated descriptor.

\section{P2-T01 interface and activation state}

The $N$ component accepts only an instantiated descriptor and returns exactly
one of the following verdicts:
\begin{itemize}[leftmargin=2em]
  \item \nolinkurl{necessary_condition_met_not_sufficient};
  \item \nolinkurl{scoped_obstruction} or
        \nolinkurl{outside_differential_scope};
  \item \nolinkurl{insufficient_physical_typing} or
        \nolinkurl{insufficient_equivalence_or_domain_data}.
\end{itemize}
It must populate source state/domain, regularity, source equivalence, output
equivalence, reconstruction, factor statistic, operational typing,
independence, stochastic scope, and escape/global scope.

This packet completes the \emph{descriptor specification}; it does not
complete a \emph{descriptor instance}. Thus $R_5=0$ and $R_6=0$ under the
preregistered activation rule. B1 remains locally frozen, B2 remains inactive,
and P4-T03 remains locked.

\section{Source-extension classification and next role}

The formal law is a new ontology primitive candidate only in the descriptive
sense that current ontology does not derive it. Its controlled statuses are
\emph{draft/control}, \emph{proposal-only}, and \emph{source-extension data}.
It is not a canonical-ontology candidate, adopted law, physical matter model,
or metric construction. Current adoption is blocked while same-milestone
continuation remains open:
\nolinkurl{blocked_adoption_open_continuation}.

Before a Candidate Constructor may instantiate the law, one separately
admitted \texttt{smuggling-auditor@0.2.0} packet should audit the descriptor,
gluing theorem, counterpair, and all 35 atomic obligations for target-atlas,
target-metric, desired-cone, benchmark-fit, metadata-authority, role-authority,
and validation-authority imports. A local audit pass would not adopt the law
or activate B2.

\section{Distance-to-GR and freeze control}

The Distance-to-GR delta is zero. The descriptor law creates no $M_{\rm src}$,
physical causal structure, signature, scale, $g_{\rm eff}$, matter coupling,
Einstein equation, Gate evidence, or benchmark status. B2 is not frozen in
this packet because the projective gluing theorem and finite-trace
underdetermination theorem are materially new mathematical payload and select
a concrete audit consequence. Repeating another generic descriptor checklist
without an audit, a constructed instance, or a precise stronger obstruction
would require freeze review.

\section{Parent-child synthesis}

The Physicist-Mathematician perspective supplied the typed ten-component law,
projective gluing theorem, finite-trace theorem, and atomic obligations. The
Physicist-Philosopher perspective separated finite protected postulate,
proposal-only continuumization, mathematical characteristic class, physical
causal interpretation, and protected adoption. Parent review resolved four
conflicts: source chart versus target atlas, schema completeness versus
instance existence, common projective zero set versus Lorentzian cone, and
scoped underdetermination versus impossibility. No conflict remains open.

\end{document}
